% Related Work section — ReCAHS paper
% Compile with natbib (\usepackage{natbib}) and sections/references.bib
% Every claim about a cited paper was checked against the PDF in papers/.

\section{Related Work}
\label{sec:related}

\paragraph{Attention-head redundancy and static pruning.}
The observation that multi-head attention is heavily over-provisioned dates
to \citet{michel2019sixteen}, who scored each head by the change in loss
when it is masked and showed that a large fraction of heads in NMT and BERT
models can be removed with little degradation. \citet{voita2019analyzing}
reached a similar conclusion with $L_0$-regularized stochastic gates,
additionally showing that the surviving heads play specialized, interpretable
roles. Subsequent work complicated this picture in two ways that our study
builds on directly. \citet{li2021differentiable} recast head pruning as a
\emph{subset-selection} problem and showed that the greedy, per-head scoring
of \citet{michel2019sixteen} systematically underestimates how many heads can
be removed, because it ignores interactions between heads.
\citet{budhraja2020weak} showed that the link between measured head
importance and actual prunability is weak: randomly chosen head subsets often
match importance-guided ones. Our results transport both critiques to
time-series forecasting and sharpen them with a controlled comparison: a
single mask chosen by \emph{joint} greedy elimination is robust across
sparsity levels, whereas the standard independent importance ranking is
fragile enough to be beaten by random pruning on minute-resolution data.

\paragraph{Dynamic and input-conditional pruning.}
A parallel line makes the pruning decision depend on the input, a special
case of conditional computation \citep{han2021dynamic}. Token-level examples
include learnable context dropping in autoregressive LLMs
\citep{anagnostidis2023dynamic} and dynamic token sparsification in vision
Transformers \citep{rao2021dynamicvit}. Closer to our setting are methods
that adapt \emph{head} usage per instance: SmartTrim prunes tokens and heads
per input in vision--language models \citep{wang2024smarttrim}, hybrid
dynamic pruning detects and skips unimportant heads at run time in
hardware \citep{jaradat2024hybrid}, delta-based reuse exploits temporal
stability of neighbouring inputs on edge devices
\citep{jelcicova2022delta}, and SPRINT selects which heads to prune
per query, showing that targeted head removal can even \emph{improve}
task performance \citep{nguyen2025structured}. These works report gains of
input-conditional sparsity over static baselines, but none of them controls
for the selection criterion: the dynamic method and the static baseline
typically differ both in \emph{how} masks are chosen and in \emph{whether}
they adapt. Our study separates these factors and finds that, for
regime-conditioned head selection in forecasting Transformers, the criterion
explains essentially all of the gain --- while an oracle over the same
per-regime masks shows that input-conditional selection retains substantial
untapped headroom, which none of five deployable routers recovers.

\paragraph{Transformers for long-term time-series forecasting.}
Our base model is PatchTST \citep{nie2023patchtst}, which applies
channel-independent patch attention and is a strong, widely used LTSF
baseline; iTransformer \citep{liu2024itransformer} instead attends across
variates, and linear models remain competitive on standard benchmarks
\citep{zeng2023dlinear}, which cautions against purely architectural claims
and motivates our efficiency framing. Structure-aware decomposition has a
long history in classical time-series analysis: STL \citep{cleveland1990stl}
separates trend, seasonal and remainder components, and regime-switching
models date back to \citet{hamilton1989regime}. Within deep forecasting,
adaptivity to temporal structure appears as architecture-level routing,
e.g.\ multi-scale adaptive pathways in Pathformer \citep{chen2024pathformer}.
We connect these threads at the level of \emph{structural sparsity}: STL
regimes define the conditioning signal, and head masks --- not architecture
--- are what adapts.

\paragraph{Pruning and retraining.}
Finally, our finding that masked-from-scratch training recovers (and on
ETTm1 exceeds) the unpruned baseline, while post-hoc fine-tuning of a
converged model fails, echoes the lottery-ticket literature
\citep{frankle2019lottery}: the value of a sparsity pattern depends on when
in training it is imposed. We did not search for winning initializations;
we only show that at $25\%$ head sparsity the pattern found on a converged
model is trainable from scratch at no accuracy cost.
